%% abstract_ko.tex — 국문초록
\begin{abstractko}

\textbf{연구 배경}: 현대 기업 환경에서는 API 문서, 인사 정책, 회의록, 보안 지침, 시스템 설계서, 출장 정책 등 다양한 형태의 지식 문서가 방대한 규모로 생성·관리되고 있다. 이러한 이종(異種) 문서 코퍼스에서 필요한 정보를 신속·정확하게 검색하는 것은 기업 운영의 핵심 과제로 부상하였다. 검색 증강 생성(Retrieval-Augmented Generation, RAG) 패러다임은 대규모 언어 모델의 생성 능력과 외부 문서 검색을 결합하여 이 문제를 해결하는 유망한 접근법으로 주목받고 있다.

\textbf{연구 목적}: 본 논문은 기업 지식 관리 시스템에서 검색 전략을 질의 특성에 따라 동적으로 선택하는 \emph{적응형 검색 계획(Adaptive Retrieval Planning)}의 효과를 실증적으로 분석한다. 구체적으로 BM25(희소 검색), 밀집 임베딩 기반 검색(Dense Retrieval), 상호 순위 융합(Hybrid RRF), 규칙 기반 계획기, GPT-5 기반 계획기 등 5가지 시스템을 체계적으로 비교하여 7개의 연구 질문(RQ1--RQ7)에 대한 실증적 답변을 제시한다.

\textbf{연구 방법}: 본 연구는 합성 기업 지식 데이터셋(Synthetic Enterprise Knowledge Dataset, SEKD)을 구축하여 통제된 실험 환경을 제공하였다. SEKD는 6개 문서 범주에 걸쳐 120개 문서, 1,206개의 섹션 인식 청크, 6가지 추론 유형(단일 홉, 다중 홉, 집계, 비교, 예외, 시간적)을 포함하는 405개의 주석 질의로 구성된다. 각 질의에는 정답 청크 ID와 전략 선택 기준(oracle) 레이블이 부여되어 정밀한 검색 평가가 가능하다. 평가 지표로는 Recall@$k$, 평균 역순위(MRR), 정규화 할인 누적 이득(nDCG@$k$), Hit Rate@$k$, 전략 선택 정확도(SSA)를 활용하였다.

\textbf{실험 결과}: 실험 결과, Hybrid RRF는 BM25 대비 MRR +15.8\%, Dense 대비 MRR +5.6\%의 성능 향상을 보여 어휘적 신호와 의미적 신호의 상보성을 입증하였다. 7개 규칙으로 구성된 결정론적 계획기는 모든 시스템 중 가장 높은 Recall@10(0.755)과 Hit@10(0.808)을 달성하였으며, 추론 비용은 \$0이었다. GPT-5 계획기는 최고의 MRR(0.533)과 nDCG@10(0.570)을 기록하였으나, 전략 선택 정확도(SSA)에서 규칙 기반 계획기(78.4\%)에 미치지 못하는 77.6\%를 기록하였고, 380건 질의 처리에 \$6.18의 API 비용이 소요되었다. 예외 처리 질의(Recall@10 $\leq$ 0.412)와 시간적 추론 질의(Dense의 경우 Recall@10 = 0.000)는 모든 시스템에서 공통적으로 낮은 성능을 보였다.

\textbf{연구 의의}: 본 연구의 주요 기여는 세 가지이다. 첫째, 적응형 검색 계획의 성능 향상 효과는 확인되었으나, 그 크기는 검색 아키텍처 선택에 의한 성능 향상(BM25→Hybrid: MRR +15.8\%)에 비해 미미한 수준(최대 MRR +2.0\%)임을 실증하였다. 둘째, 결정론적 규칙 기반 계획기가 GPT-5와 동등한 수준의 검색 성능을 수백만 배 낮은 지연 시간과 \$0 비용으로 달성함을 보였다. 셋째, 예외 처리 및 시간적 추론 질의에 대한 구조적 취약점을 규명하고, 이를 해결하기 위한 전문화된 검색 메커니즘의 필요성을 제시하였다. 이러한 결과는 기업 RAG 시스템 설계 시 검색 아키텍처 선택을 최우선 과제로 삼아야 한다는 실용적 지침을 제공한다.

\bigskip
\noindent\textbf{핵심어}: 검색 증강 생성(RAG), 기업 지식 관리, 적응형 검색, BM25, 밀집 검색, 하이브리드 검색, 상호 순위 융합(RRF), 대규모 언어 모델, 검색 계획, 정보 검색 평가

\end{abstractko}
